\documentclass[11pt,a4paper]{article}
% Kodowanie i obsługa języka polskiego
\usepackage[utf8]{inputenc}      % Kodowanie wejścia
\usepackage[T1]{fontenc}         % Kodowanie fontów
\usepackage[polish]{babel}       % Obsługa języka polskiego

% Pakiety matematyczne i inne przydatne
\usepackage{mathtools}
\usepackage{tabularx}
\usepackage{array} % for centering in p-columns
\usepackage{calc}
\usepackage[makeroom]{cancel}
\usepackage{xcolor}
\usepackage{bm}
\usepackage{marginnote}
\usepackage{tikz}
\usetikzlibrary{trees}
\usetikzlibrary{decorations.pathreplacing}
\usepackage{amsmath, amssymb, amsthm}  % Pakiety do matematyki
\usepackage{graphicx}                  % Obsługa grafiki
\usepackage{hyperref}                  % Linki i spis treści
\usepackage{geometry}                  % Ustawienie marginesów
\usepackage{algorithm}               % Algorytmy
\usepackage{algpseudocode}           % Pseudokod
\usetikzlibrary{positioning, shapes.geometric, arrows.meta}
\geometry{margin=2cm}                  % Ustawienie marginesów na 2 cm
\hypersetup{
    colorlinks=true,
    linkcolor=blue,
    filecolor=magenta,
    urlcolor=cyan,
}
\title{Notatki z Kryptografii}
\author{Jakub Kogut}

\begin{document}

\maketitle

\tableofcontents  % Spis treści (opcjonalnie)
\newpage
\section{Wstęp do Wykładów}
Zaliczenie przez ćwiczenia, prawdopodobnie będą 2 kolokwia/kartkówki, jak i również listy zadań. 
\paragraph{Kilka słów wstępu}
Generalne operacje kryptograficzne wykonuje się na elementach pewnych grup. Często wykożystyje się do tego grupy multiplikatywne $\mathbb{Z}_p^*$, gdzie $p$ jest liczbą pierwszą. Wykorzystuje się również grupy eliptyczne, które są bardziej skomplikowane, ale oferują większe bezpieczeństwo przy mniejszych rozmiarach kluczy.

\end{document}
